% elimtable.tex -- generated by reviews/elimtable-workdir/elimtable.py
% Needs only booktabs (\toprule,\midrule,\bottomrule) and array, both
% already loaded by PAPER-n6-872.tex.  No longtable.  Uses the paper's
% macros \slack, \Ncyc, \NCh and amsmath's \text.
%
% Suggested accompanying paragraph (to sit next to the table):
%   Write B = D/5+X4+Z = W-142.  A type of index B is admissible exactly
%   at the lengths 867+B <= L <= 871, with slack = L-867-B <= 3.  R4 is
%   the only slack-dependent rejection rule and floor(slack/4)=0
%   throughout, so a type's verdict does not depend on L; this was
%   checked for every admissible slack.  For an empty type the rules
%   named are a minimal sufficient set: R1-R3 and R7 (the derived-count
%   non-negativities and the defect split) together with the listed
%   rules already empty the type, and deleting any listed rule leaves a
%   survivor.  The boldface rule is the deciding one and the last column
%   gives its two sides on a profile that satisfies every other listed
%   rule, at the largest e for which such a profile exists.

\begin{table}[htbp]\centering\footnotesize
\setlength{\tabcolsep}{4pt}
\caption{Elimination record behind Lemma~\ref{lem:870} ($L=868,869$), Proposition~\ref{prop:871} ($L=870$) and Theorem~\ref{thm:classes} ($L=871$), one row per structural type.  $B=D/5+X4+Z=W-142$; a type of index $B$ occurs exactly at the lengths $867+B\le L\le871$ and its verdict is independent of $\slack$.  The $e$-column is the range left by R1.  For an empty type the named rules, with R1--R3 and R7, already empty it, and deleting any one of them leaves a survivor; the boldface rule is the deciding one and the last column is its two sides on a profile satisfying all the others.  Deleting the two surviving classes with $e=0$ (Lemma~\ref{lem:e0}) leaves the sixteen of Theorem~\ref{thm:classes}.}
\label{tab:elim}
\begin{tabular}{@{}clcl>{\raggedright\arraybackslash}p{29mm}>{\raggedright\arraybackslash}p{51mm}@{}}
\toprule
$B$ & type & $e$ & verdict & rules & deciding numbers\\
\midrule
1 & $(0,0,1)$ & $\le24$ & empty & R4+\allowbreak R8a+\allowbreak \textbf{R10} & $F_p+S_5=24>4R_{\mathrm{cap}}=4$ \ ($e=1$)\\
 & $(0,1,0)$ & $\le23$ & empty & R4+\allowbreak R6+\allowbreak \textbf{R10} & $F_p+S_5=24>4R_{\mathrm{cap}}=8$ \ ($e=0$)\\
 & $(5,0,0)$ & $\le24$ & empty & R4+\allowbreak R9a+\allowbreak \textbf{R12} & $P-2D=15>\text{R12 cap}=4$ \ ($e=24$)\\
\midrule
2 & $(0,0,2)$ & $\le25$ & empty & R4+\allowbreak R8a+\allowbreak \textbf{R10} & $F_p+S_5=24>4R_{\mathrm{cap}}=4$ \ ($e=2$)\\
 & $(0,1,1)$ & $\le24$ & empty & R4+\allowbreak R8a+\allowbreak \textbf{R10} & $F_p+S_5=24>4R_{\mathrm{cap}}=8$ \ ($e=1$)\\
 & $(0,2,0)$ & $\le23$ & empty & R4+\allowbreak R6+\allowbreak \textbf{R10} & $F_p+S_5=24>4R_{\mathrm{cap}}=12$ \ ($e=0$)\\
 & $(5,0,1)$ & $\le25$ & empty & R4+\allowbreak R9a+\allowbreak \textbf{R12} & $P-2D=15>\text{R12 cap}=4$ \ ($e=25$)\\
 & $(5,1,0)$ & $\le24$ & empty & R4+\allowbreak R9a+\allowbreak \textbf{R12} & $P-2D=15>\text{R12 cap}=4$ \ ($e=24$)\\
 & $(10,0,0)$ & $\le25$ & empty & R4+\allowbreak R9a+\allowbreak \textbf{R12} & $P-2D=6>\text{R12 cap}=4$ \ ($e=25$)\\
\midrule
3 & $(0,0,3)$ & $\le26$ & empty & R4+\allowbreak R8a+\allowbreak \textbf{R10} & $F_p+S_5=24>4R_{\mathrm{cap}}=4$ \ ($e=3$)\\
 & $(0,1,2)$ & $\le25$ & empty & R4+\allowbreak R8a+\allowbreak \textbf{R10} & $F_p+S_5=24>4R_{\mathrm{cap}}=8$ \ ($e=2$)\\
 & $(0,2,1)$ & $\le24$ & empty & R4+\allowbreak R8a+\allowbreak \textbf{R10} & $F_p+S_5=24>4R_{\mathrm{cap}}=12$ \ ($e=1$)\\
 & $(0,3,0)$ & $\le23$ & empty & R4+\allowbreak R6+\allowbreak \textbf{R10} & $F_p+S_5=24>4R_{\mathrm{cap}}=16$ \ ($e=0$)\\
 & $(5,0,2)$ & $\le26$ & empty & R4+\allowbreak R9a+\allowbreak \textbf{R12} & $P-2D=15>\text{R12 cap}=4$ \ ($e=26$)\\
 & $(5,1,1)$ & $\le25$ & empty & R4+\allowbreak R9a+\allowbreak \textbf{R12} & $P-2D=15>\text{R12 cap}=4$ \ ($e=25$)\\
 & $(5,2,0)$ & $\le24$ & empty & R4+\allowbreak R9a+\allowbreak \textbf{R12} & $P-2D=15>\text{R12 cap}=4$ \ ($e=24$)\\
 & $(10,0,1)$ & $\le26$ & empty & R4+\allowbreak R9a+\allowbreak \textbf{R13c} & $P=26>F(26,10)$: unrealisable \ ($e=26$)\\
 & $(10,1,0)$ & $\le25$ & empty & R4+\allowbreak R6+\allowbreak R9a+\allowbreak R11b+\allowbreak \textbf{R13} & $P=26>F_s(\NCh{=}12,D{=}10,\mathit{PEN}{=}10)=18$ \ ($e=10$)\\
 & $(15,0,0)$ & $\le26$ & empty & R4+\allowbreak R9a+\allowbreak R11b+\allowbreak R14+\allowbreak \textbf{R13r} & $P=27>F_r(\NCh{=}16,S_5{=}15,S_4{=}0)=19$ \ ($e=15$)\\
\midrule
4 & $(0,0,4)$ & $\le27$ & empty & R4+\allowbreak R8a+\allowbreak \textbf{R10} & $F_p+S_5=24>4R_{\mathrm{cap}}=4$ \ ($e=4$)\\
 & $(0,1,3)$ & $\le26$ & empty & R4+\allowbreak R8a+\allowbreak \textbf{R10} & $F_p+S_5=24>4R_{\mathrm{cap}}=8$ \ ($e=3$)\\
 & $(0,2,2)$ & $\le25$ & empty & R4+\allowbreak R8a+\allowbreak \textbf{R10} & $F_p+S_5=24>4R_{\mathrm{cap}}=12$ \ ($e=2$)\\
 & $(0,3,1)$ & $\le24$ & empty & R4+\allowbreak R8a+\allowbreak \textbf{R10} & $F_p+S_5=24>4R_{\mathrm{cap}}=16$ \ ($e=1$)\\
 & $(0,4,0)$ & $\le23$ & empty & R4+\allowbreak R6+\allowbreak \textbf{R10} & $F_p+S_5=24>4R_{\mathrm{cap}}=20$ \ ($e=0$)\\
 & $(5,0,3)$ & $\le27$ & empty & R4+\allowbreak R9a+\allowbreak \textbf{R12} & $P-2D=15>\text{R12 cap}=4$ \ ($e=27$)\\
 & $(5,1,2)$ & $\le26$ & empty & R4+\allowbreak R9a+\allowbreak \textbf{R12} & $P-2D=15>\text{R12 cap}=4$ \ ($e=26$)\\
 & $(5,2,1)$ & $\le25$ & empty & R4+\allowbreak R9a+\allowbreak \textbf{R12} & $P-2D=15>\text{R12 cap}=4$ \ ($e=25$)\\
 & $(5,3,0)$ & $\le24$ & empty & R4+\allowbreak R6+\allowbreak R9a+\allowbreak R11b+\allowbreak R12+\allowbreak \textbf{R13c} & $P=25>F(\NCh{=}5,D{=}5)=23$ \ ($e=1$)\\
 & $(10,0,2)$ & $\le27$ & empty & R4+\allowbreak R9a+\allowbreak R11b+\allowbreak R13+\allowbreak \textbf{R14} & $26$ link-paths do not pack into $N_{\mathrm{comp}}{=}21$ groups \ ($e=7$)\\
 & $(10,1,1)$ & $\le26$ & empty & R4+\allowbreak R9a+\allowbreak R11b+\allowbreak R14+\allowbreak \textbf{R13r} & $P=26>F_r(\NCh{=}12,S_5{=}10,S_4{=}0)=18$ \ ($e=11$)\\
 & $(10,2,0)$ & $\le25$ & \textbf{alive} & --- & $e=0,2$ \ (2 classes)\\
 & $(15,0,1)$ & $\le27$ & \textbf{alive} & --- & $e=2$ \ (1 class)\\
 & $(15,1,0)$ & $\le26$ & \textbf{alive} & --- & $e=0,\dots,4,7$ \ (6 classes)\\
 & $(20,0,0)$ & $\le27$ & \textbf{alive} & --- & $e=1,\dots,9$ \ (9 classes)\\
\bottomrule
\end{tabular}
\end{table}

% --- the sixteen classes, read straight off the $B=4$ block above ---
% (10,2,0) e=2 ; (15,0,1) e=2 ; (15,1,0) e=1,2,3,4,7 ; (20,0,0) e=1..9
